%klassen  ------------------------
\question[9]
a. Implementiere eine Klasse \texttt{Filmmusik} mit den Attributen
\texttt{film} und \texttt{komponist}, so dass
die folgenden Filmmusik-Objekte erzeugt werden können:
\begin{lstlisting}
m1 = Filmmusik("Jurassic Park","John Williams")
m2 = Filmmusik("Star Wars","John Williams")
m3 = Filmmusik("Terminal","John Williams")
m4 = Filmmusik("Inception","Hans Zimmer")
m5 = Filmmusik("Interstellar","Hans Zimmer")
m6 = Filmmusik("Lord of the Rings","Howard Shore")
\end{lstlisting}
b. Implementiere eine string-Methode, die mit \texttt{print(m1)}
folgende Ausgabe erzeugt: \\
\texttt{Jurassic Park - Musik: John Williams} \\
c. Erzeuge eine Liste \texttt{fmlist} mit den Objekten m1, m2... m6.
Die Objekte m1 bis m6 können als erzeugt vorausgesetzt werden. \\
d. Gehe mit einer Schleife durch die \texttt{fmlist} und gib
alle Filmtitel mit Musik von John Williams aus. \\
e. Erzeuge für die 6 Filme in \texttt{fmlist} ein Dictionary
\texttt{fmmap}, bei dem jedem Komponistennamen (key) eine Liste seiner
Filmmusik-Objekte als value zugeordnet ist. \\

Beispiel einer Nutzung von \texttt{fmmap}:
\begin{lstlisting}
>>> for f in fmmap["Hans Zimmer"]:
  	print(f)

Inception - Musik: Hans Zimmer
Interstellar - Musik: Hans Zimmer
\end{lstlisting}

\begin{solutionbox}{7.5cm}
\begin{lstlisting}
class Filmmusik:
    def __init__(self,film,komponist):
        self.film = film
        self.komponist = komponist

    def __str__(self):
        return self.film + " - Musik: " + self.komponist


fmlist = [m1,m2,m3,m4,m5,m6]
for f in fmlist:
    if f.komponist == "John Williams":
        print(f.film)

fmmap =  {"John Williams":[m1,m2,m3],
              "Hans Zimmer":[m4,m5],"Howard Shore":[m6]}
\end{lstlisting}
\end{solutionbox}
